\section{The formalization}
\label{sec:formalization}

The Lean development is the directory \lean{ImprovedExponents} of the
repository \url{https://github.com/qawbecrdtey/ImprovedExponents}, which also
contains this paper.  It has \LeanFiles{} files and
\LeanLines{} lines, of which \LeanCellLines{} are generated certificates for
Theorems~\ref{thm:global}, \ref{thm:global-full} and~\ref{prop:shapes}; about
\LeanHandLines{} lines are written by hand.

\subsection{Relation to the upstream formalization}

The part of the formalization \cite{upstream} of the paper that we use, the
\UpstreamFiles{} files of Lean that our development imports directly or
transitively, is included in the repository at a fixed revision
(\texttt{\UpstreamRev}; directory \texttt{upstream/3sum-apsp}) and ported
from Lean and Mathlib \texttt{\UpstreamVersion}, for which it was written, to
\texttt{\LeanVersion}.  The port changes only what no longer compiles, in
\UpstreamChanged{} of the files, and no statement or definition; every other
file is upstream's, byte for byte.  It provides the word RAM
and the notion \lean{Problem.SolvedInTime} (file \lean{EndStatement.lean}), a
small programming language with a time model (\lean{lightModel}) and a compiler
to the word RAM, the programs of the paper, and the proofs of its theorems.

We reuse the following.
\begin{itemize}
\item The program of Section~4, \lean{program31}, for arbitrary rational
  parameters, with its correctness and the bound of its running time in terms
  of the two costs of Theorem~30.
\item The count of the boxes, \lean{lemma_29_count}, which leaves the tail
  \(\sum_{d\ge t}\beta_d\) explicit.
\item The host of Theorem~17 and the passage from a host to a claim about
  programs, \lean{claim17_of_host}.
\item The reductions of Theorem~21 and the passage to the end statements:
  \lean{threeSum_of_uniform}, \lean{minPlus_polylog_of_uniform},
  \lean{apsp_polylog_of_uniform} and \lean{SolvedInTime.endStatement}.
\item The enclosures of logarithms between rational numbers, \lean{log_mem}.
\end{itemize}
Upstream proves the bounds of the paper at the parameters of the paper.  Where
an upstream proof fixes parameters, we repeat it for general ones: the solver
of the thin product at arbitrary ratios (\lean{Pipeline/ThinClaim.lean}), the
time of the program against the exact tail (\lean{Cost8X}), the routines that
compute the parameters of the host (\lean{ParamRoutines}), and the arithmetic
of the total time (\lean{Total}).  Where an upstream proof assumes the full
encodings, we repeat it under the weaker contract of
Section~\ref{sec:encoder-program}: the encoder of Theorem~5, the procedures of
Theorem~30 and the routines with chosen parameters (\lean{PrunedProgram}), and
the time of the new program against \lean{cost8P} (\lean{Cost8P}).  The file
\texttt{NOTICE} of the repository lists the files that are adapted from
upstream, which is distributed under the Apache~2.0 license.

\subsection{What is new}

Table~\ref{tab:dirs} lists the directories.

\begin{table}[ht]
\centering\small
\begin{tabular}{lr>{\raggedright\arraybackslash}p{3.3in}}
\toprule
directory & lines & content\\
\midrule
\lean{ExactCount} & \LeanExactCountLines & \(\gtwo\), Theorem~\ref{thm:tail},
  Proposition~\ref{prop:lemmaA}\\
\lean{Cost8X} & \LeanCostEightXLines & the upstream program against the exact tail
  (Theorem~\ref{thm:offline})\\
\lean{PrunedEncoding} & \LeanPrunedEncodingLines & Theorem~\ref{thm:pruned},
  Proposition~\ref{prop:costP}\\
\lean{PrunedProgram} & \LeanPrunedProgramLines & the pruned encoder and the
  solver under the weaker contract (Section~\ref{sec:encoder-program})\\
\lean{Cost8P} & \LeanCostEightPLines & the new program against \lean{cost8P}
  (Theorem~\ref{thm:encoder})\\
\lean{HostMid} & \LeanHostMidLines & the host with one added assignment
  (Theorem~\ref{thm:mid})\\
\lean{ParamRoutines} & \LeanParamRoutinesLines & routines for \(\lfloor n^{a/b}\rfloor\) and
  \(\lceil D^{c/d}\rceil\)\\
\lean{Pipeline} & \LeanPipelineLines & thin-product claims, the regime
  tests, the offline routines, the general theorems\\
\lean{Total} & \LeanTotalLines & Theorem~\ref{thm:total}\\
\lean{Optimum} & \LeanOptimumLines & Section~\ref{sec:optimum},
  Proposition~\ref{prop:half}, the shapes of Section~\ref{sec:limits};
  \LeanCellLines{} lines are generated\\
\lean{MinPlus} & \LeanMinPlusLines & Lemma~\ref{lem:F},
  Proposition~\ref{prop:minplus-math}\\
\lean{AllEdges} & \LeanAllEdgesLines & the all-edges route to APSP
  (Theorems~\ref{thm:pairs}, \ref{thm:mid-ae}, \ref{thm:apsp-ae})\\
\lean{Statements.lean} & \LeanStatementsLines & the theorems of
  Section~\ref{sec:consequences}\\
\bottomrule
\end{tabular}
\caption{The Lean development.}
\label{tab:dirs}
\end{table}

Two points of the formalization differ from a proof on paper.

\emph{The programs are almost fixed.}  In Theorems~\ref{thm:T2}
and~\ref{thm:T3} the program for the thin product and the host of the
reduction are upstream's, at other parameters; the only new program text is
that of the routines that compute \(D'=\lfloor n^{a/b}\rfloor\) and
\(g'=\lceil D'^{c/d}\rceil\).  So the improvement from \(\ETpaper\) to
\(\dThree\) is obtained by a new analysis of existing programs.
Theorem~\ref{thm:T4} adds one assignment to the host and three short
procedures for the regime test; the all-edges route (Theorem~\ref{thm:apsp-ae})
adds the two hosts of Section~\ref{sec:alledges}, which reuse upstream's
passes, scans, hashing and instance construction.  Theorem~\ref{thm:T1} adds
seven procedures (\lean{programP}): the pruned encoder and six that call it,
or call upstream's procedures under the new contract; the procedures that
build the tiles and answer the queries are upstream's, with a new proof.

\emph{Numerals are corollaries.}  The general theorems
(\lean{allSolved_paper}, \lean{allSolved_exact}, \lean{allSolved_full},
\lean{allSolved_pruned}) quantify over the parameters or over all savings
below the closed form \(S(c)\).  The decimal exponents of
Table~\ref{tab:results} follow from certified enclosures of \(S(c)\) at one
ratio (Proposition~\ref{prop:values}).

\subsection{The main statements}

Table~\ref{tab:names} lists the Lean names of the main statements.

\begin{table}[ht]
\centering\footnotesize
\setlength{\tabcolsep}{4pt}
\begin{tabular}{lll}
\toprule
statement & declaration & file\\
\midrule
Theorem~\ref{thm:tail} & \lean{sum_beta_le_exp} & \lean{ExactCount/Tail.lean}\\
Proposition~\ref{prop:lemmaA} & \lean{sum_card_Leaves_le_exp}
  & \lean{ExactCount/Leaves.lean}\\
Theorem~\ref{thm:offline} & \lean{offlineWithinX} & \lean{Cost8X/Offline.lean}\\
Theorem~\ref{thm:pruned} & \lean{pruned_encoding}
  & \lean{PrunedEncoding/Pruned.lean}\\
Proposition~\ref{prop:costP} & \lean{costsP} & \lean{PrunedEncoding/Costs.lean}\\
Remark~\ref{rem:encwork} & \lean{encWorkP_le_M_div} & \lean{PrunedProgram/Work.lean}\\
Theorem~\ref{thm:encoder} & \lean{prunedEncoderClaim}
  & \lean{Pipeline/PrunedClaim.lean}\\
 & \lean{encodeP_spec}, \lean{tEncP_le}
  & \lean{PrunedProgram/Encode/Recursion.lean}\\
 & \lean{offline32P_meets} & \lean{PrunedProgram/Wrapper/Offline.lean}\\
 & \lean{offlineWithinP} & \lean{Cost8P/Offline.lean}\\
Theorem~\ref{thm:mid} & \lean{claim17Mid_of_paramProcs}
  & \lean{HostMid/AtParameters.lean}\\
Theorem~\ref{thm:total} & \lean{explicitFrom_mid} & \lean{Total/Explicit.lean}\\
Theorem~\ref{thm:sup} & \lean{isLUB_tupleSaving}
  & \lean{Optimum/PerCSupremum.lean}\\
Proposition~\ref{prop:rational} & \lean{exists_rat_tuple_basePruned}
  & \lean{Optimum/PerCDensity.lean}\\
Proposition~\ref{prop:half} & \lean{isLUB_paperSaving},
  \lean{savingF_GamHalf_lt} & \lean{Optimum/PaperAccounting.lean}\\
Theorem~\ref{thm:global} & \lean{optimum_global} & \lean{Statements.lean}\\
Theorem~\ref{thm:global-full} & \lean{optimum_global_full} & \lean{Statements.lean}\\
Theorem~\ref{thm:T2} & \lean{threeSum_paper}, \lean{apsp_paper}
  & \lean{Statements.lean}\\
Theorem~\ref{thm:T3} & \lean{threeSum_exact}, \lean{apsp_exact}
  & \lean{Statements.lean}\\
Theorem~\ref{thm:T4} & \lean{allSolved_full}, \lean{threeSum_full}
  & \lean{Statements.lean}\\
Theorem~\ref{thm:T1} & \lean{allSolved_pruned} & \lean{Pipeline/PrunedClaim.lean}\\
 & \lean{threeSum_pruned}, \lean{apsp_pruned} & \lean{Statements.lean}\\
Lemma~\ref{lem:F} & \lean{add_lt_iff_shiftDiff}
  & \lean{MinPlus/Comparison.lean}\\
Proposition~\ref{prop:minplus-math} & \lean{search_exactTriangle}
  & \lean{MinPlus/ExactTriangle.lean}\\
Theorem~\ref{thm:pairs} & \lean{isHost_pairsAE} & \lean{AllEdges/Pairs/Host.lean}\\
Theorem~\ref{thm:mid-ae} & \lean{ae17_solves} & \lean{AllEdges/Host/ProgramAE.lean}\\
 & \lean{midHostRat_ae} & \lean{AllEdges/Host/ClaimAE.lean}\\
Theorem~\ref{thm:apsp-ae} & \lean{minPlusSolved_full}, \lean{apsp_allEdges}
  & \lean{Statements.lean}\\
 & \lean{minPlusSolved_pruned}, \lean{apsp_allEdges_pruned}
  & \lean{Statements.lean}\\
Proposition~\ref{prop:shapes} & \lean{shapes_table} & \lean{Optimum/Shapes/Table.lean}\\
\bottomrule
\end{tabular}
\caption{The main statements and their Lean names.  The end statements for
Exact Triangle, the \(\minplus\)-product and APSP are named in the same way
(\texttt{exactTriangle\_full}, \texttt{minPlus\_full}, \texttt{apsp\_full},
\texttt{minPlus\_allEdges}, and so on).}
\label{tab:names}
\end{table}

\subsection{What has to be trusted}

For the end statements of Theorems~\ref{thm:T2}, \ref{thm:T3}, \ref{thm:T4},
\ref{thm:T1} and~\ref{thm:apsp-ae}: the Lean kernel, and the definitions of
upstream's \lean{EndStatement.lean}, which state what it means that a problem is
solved in time \(O(n^{r})\) on the word RAM.  That file imports nothing, and
the module \texttt{ImprovedChallenge.lean} of the repository, which states the
four strongest end statements, contains its 132 lines of Lean verbatim and
imports nothing either.  An end statement such as
\begin{center}
\leanraw{theorem threeSum_pruned : EndStatement.ThreeSum.SolvedInTime 1.99896}
\end{center}
mentions no definition of ours and takes no argument.  No statement of this
paper rests on a hypothesis.

The development contains no \texttt{sorry} and declares no axiom.  Lean
reports that the end statements depend on the axioms \texttt{propext},
\texttt{Classical.choice} and \texttt{Quot.sound}, the three axioms that
Mathlib uses throughout, and on no other; in particular the numerical
certificates are checked by the kernel and not by compiled code.

\subsection{How to check}

In the root of the repository, \texttt{lake exe cache get} fetches the compiled
Mathlib, and \texttt{lake build} checks everything: the included part of
upstream, the library and the two modules described next; the script
\texttt{scripts/check-upstream.sh} compares the included part with upstream's
revision and confirms that every changed file is marked and listed.  The module \texttt{ImprovedChallenge.lean}
states the four strongest end statements, \lean{threeSum_pruned}
(\(\ThreeSumOne\)), \lean{exactTriangle_pruned} (\(\ETOne\)),
\lean{minPlus_allEdges_pruned} and \lean{apsp_allEdges_pruned}
(\(\APSPFnOne\); Theorems~\ref{thm:T1} and~\ref{thm:apsp-ae}), with
upstream's \lean{EndStatement.lean} copied into it verbatim (the script
\texttt{scripts/check-endstatement.sh} confirms the copy against the included
file, which is upstream's) and nothing imported, each with \texttt{sorry} for
its proof; the module \texttt{ImprovedSolution.lean} imports the proofs.  The
two modules and the configuration \texttt{comparator.json} are in the form
that the tool Comparator \cite{comparator} checks: that the solution proves
the statements of the challenge, literally, that it uses only the axioms
\texttt{propext}, \texttt{Quot.sound} and \texttt{Classical.choice}, and that
everything replays in Lean's kernel and in independent kernels; the
compilations run in a sandbox.  From Lean \texttt{v4.35.0-rc2} on, Comparator
ships with the toolchain as \texttt{lake comparator}, with the independent
kernels NanoDa and con-ron.  The script \texttt{scripts/verify-comparator.sh}
runs it as the Palomar registry does, in a bubblewrap sandbox and in a fresh
copy of the sources without build products, so that the challenge, the
solution and everything they import are compiled inside the sandbox.  We ran it
on the final development: \texttt{lake comparator} accepted the pair on
\ComparatorDate{} (``Your solution is okay!'', all three kernels; the log is
\texttt{scripts/comparator-run.log}).  Every Lean
file of the repository uses Lean's module system, and the statement module
imports nothing, as the Palomar registry \cite{Palomar} requires of a
submission.

The certificates of Theorems~\ref{thm:global} and~\ref{thm:global-full} are
regenerated by \texttt{python3 -m search.certs}, those of
Proposition~\ref{prop:shapes} with the option \texttt{--encoding shapes}, and
the script
\texttt{check\_lean\_names.py} next to the source of this paper checks that
every Lean name cited here exists in the sources.

\subsection{What is not formalized}

Only parts of Section~\ref{sec:limits}: the discussion of mixtures of
shapes, the grid of hypothetical shapes, the experiments on the arrangement of
the wanted set, and the remark on \(\omega\).  Each is marked ``\nf'' there.
Everything else, including the comparison of the shapes of Sch\"onhage's
family (Proposition~\ref{prop:shapes}), is machine-checked.
